% !TEX program = xelatex
\documentclass{英语试卷}
\begin{document}
\examtitle{金华十校 2025 年 4 月高三模拟考试}{英语}
\par 金华十校 2025 年 4 月高三模拟考试\par
\par 英语学科 试题卷\par
\begin{notice}
\par 1. 答卷前，考生务必将自己的姓名、准考证号填写在答题卡上。\par
\par 2. 回答选择题时，选出每小题答案后，用铅笔把答题卡上对应题目的答案标号涂黑。如需改动，用橡皮擦干净后，再选涂其他答案标号。回答非选择题时，将答案写在答题卡上，写在本试卷上无效。\par
\par 3. 考试结束后，将本试卷和答题卡一并交回。\par
\end{notice}
\parttitle{第一部分 听力（共两节，满分 30 分）}
\par 做题时，先将答案标在试卷上，录音内容结束后，你将有两分钟的时间将试卷上的答案转涂到答题卡上。\par
\sectitle{第一节（共5小题：每小题1.5分，满分7.5分）}
\par 听下面 5 段对话。每段对话后有一个小题，从题中所给的 A、B、C 三个选项中选出最佳选项。听完每段对话后，你都有 10 秒钟的时间来回答有关小题和阅读下一小题。每段对话仅读一遍。\par
\question{1. How does the woman feel?}
\begin{choices}(3)
  \item Content.
  \item Disturbed.
  \item Stressed.
\end{choices}
\question{2. Where are the speakers probably?}
\begin{choices}(3)
  \item At school.
  \item In a forest.
  \item At a news agency.
\end{choices}
\question{3. What's the probable relationship between the speakers?}
\begin{choices}(3)
  \item Friends.
  \item Husband and wife.
  \item Employer and employee.
\end{choices}
\question{4. What does the man like about chopsticks?}
\begin{choices}(3)
  \item Their usefulness.
  \item Their appearance.
  \item Their convenience.
\end{choices}
\question{5. What does the man want to do?}
\begin{choices}(2)
  \item Cancel an appointment. 3. Book a repair service.
  \item Postpone a schedule.
\end{choices}
\sectitle{第二节（共15小题：每小题1.5分，满分22.5分）}
\par 听下面 5 段对话或独白。每段对话或独白后有几个小题，从题中所给的 A、B、C 三个选项中选出最佳选项。听每段对话或独白前，你将有时间阅读各个小题，每小题 5 秒钟；听完后，各小题将给出 5 秒钟的作答时间。每段对话或独白读两遍。\par
\question{6. What is the main purpose of the speakers' hike?}
\begin{choices}(3)
  \item To refresh themselves.
  \item To get some exercise.
  \item To explore a lake.
\end{choices}
\question{7. How far do the speakers still have to go to reach the lake?}
\begin{choices}(3)
  \item 4.5 miles.
  \item 6 miles.
  \item 7.5 miles.
\end{choices}
\par 听第 7 段材料，回答第 8 至 10 题。\par
\question{8. What do we know about the woman?}
\begin{choices}(3)
  \item She is an animal lover.
  \item She never spoils her pet.
  \item She keeps a bird as a pet.
\end{choices}
\question{9. What is Joey?}
\begin{choices}(3)
  \item A bird.
  \item A cat.
  \item A dog.
\end{choices}
\question{10. What's wrong with Kevin?}
\begin{choices}(3)
  \item He is sick.
  \item He is sad.
  \item He is heavy.
\end{choices}
\par 听第 8 段材料，回答第 11 至 13 题，\par
\question{11. What are the speakers doing?}
\begin{choices}(3)
  \item Cooking dinner.
  \item Having a meal.
  \item Watching a show.
\end{choices}
\question{12. What is important to understand Peking Opera according to the woman?}
\begin{choices}(3)
  \item Understanding the patterns of facial makeup.
  \item Having a basic knowledge of the Chinese language
  \item Learning about the stories behind those historical characters.
\end{choices}
\question{13. What will the woman do next?}
\begin{choices}(3)
  \item Go to her company.
  \item Catch a bus.
  \item Book a taxi.
\end{choices}
\par 听第 9 段材料，回答第 14 至 17 题。\par
\question{14. What does the woman ask about the man's travel?}
\begin{choices}(3)
  \item Who he will go with.
  \item When he plans to leave.
  \item Why he wants to take the trip
\end{choices}
\question{15. When does the man plan to return?}
\begin{choices}(3)
  \item At the beginning of March.
  \item In the middle of March.
  \item At the end of March.
\end{choices}
\question{16. What does the man ask his neighbor to take care of?}
\begin{choices}(3)
  \item His dog.
  \item His car.
  \item His house.
\end{choices}
\question{17. What does the man do?}
\begin{choices}(3)
  \item He is a bank clerk.
  \item He is an engineer.
  \item He is a financial advisor.
\end{choices}
\par 听第 10 段材料，回答第 18 至 20 题。\par
\question{18. What is the main idea of the talk?}
\begin{choices}(3)
  \item Why to travel by train.
  \item What to do on a long train trip.
  \item How to take care of yourself on a train.
\end{choices}
\question{19. Why is reading more pleasant on a train than on a bus?}
\begin{choices}(3)
  \item The scenery is more beautiful.
  \item The seats are more comfortable.
  \item You are less likely to feel dizzy.
\end{choices}
\question{20. What does the speaker suggest travelers do to stay occupied on the train?}
\begin{choices}(3)
  \item Make phone calls.
  \item Learn a new skill.
  \item Have an online meeting.
\end{choices}
\parttitle{第二部分 阅读理解（共两节，满分 50 分）}
\sectitle{第一节（共15小题：每小题2.5分，满分37.5分）}
\par 阅读下列短文，从每题所给的 A、B、C 和 D 四个选项中，选出最佳答案。\par
\passagelabel{A}
\par A Poster Contest\par
\par Engaging people in fire safety and prevention from an early age is crucial, so the State Fire Marshal's Office is inviting primary school students across Ohio to participate in its Fire Safety Poster Contest. The contest is held annually to engage young minds, emphasizing various issues related to fire safety education. Contest rules\par
\par Posters must include a fire prevention theme\par
\par Posters must display positive behavior statements of what students can do to prevent a fire at home.\par
\par Posters must only be on 8.5 x 11-inch paper.\par
\par Any colors may be used; however, a white or light-colored background is preferred.\par
\par Crayons, markers, pencils and similar tools can be used.\par
\par All artwork must be original. Computer-generated images and clippings from magazines and newspapers will not be accepted.\par
\par The poster surface should remain flat.\par
\par Contest procedures\par
\par Elementary schools across the state must turn over the posters to their local fire departments prior to the September 30 deadline, so those entries can be displayed during Fire Prevention Week (October 6-12). Then each of the local departments must select two winning posters — one from each grade category (grades 1-3 and 4-6) and submit them to the State Fire Marshal's Office by November 15. A committee will review all the submitted posters and select six winners in each grade category. Winning entries will be announced next spring.\par
\par Additional information on the contest is available at com.ohio.gov/SFMPosterContest.\par
\question{21. Why are primary school students encouraged to join in the contest?}
\begin{choices}(2)
  \item To stimulate their creative minds.
  \item To display their brilliant artistic skills.
  \item To learn more about local fire departments.
  \item To raise their awareness of fire safety and prevention.
\end{choices}
\question{22. Which of the following will probably result in disqualification?}
\begin{choices}(2)
  \item Drawing on a 8.5 x 11-inch paper.
  \item Using AI tools to generate images
  \item Centering on a fire prevention theme
  \item Choosing a dark-colored background.
\end{choices}
\question{23. How many winning entries will be announced next spring?}
\begin{choices}(2)
  \item 2.
  \item 4.
  \item 6.
  \item 12.
\end{choices}
\passagelabel{B}
\par Jared was my regular hiking, ice-climbing and mountain-biking fellow. We were on an overnight trip, just a single night, and what he'd brought for dinner was this: Rice. Literally nothing else.\par
\par This was not a result of him trying to pack as light as possible, nor was it because he was broke. He wasn't. I tried to convince him to at least put soy sauce(酱油) on it. But no, Jared wasn't having a bar of it. "I don't need it," he said flatly.\par
\par Now, if he was talking only about outdoors' meals, I'd almost get it. When you're ravenous after a day of burning off calories walking or climbing or backcountry skiing in the mountains, how awesome do things taste! I know what you're thinking; it's simply that, as the saying goes, hunger is the best sauce. And I agree. Absolutely.\par
\par But there's more to it than that. Last month, I went down the coast car camping with my family. One morning, my wife made a statement. "I don't see the point," she said. "It's so much work, camping. But the views at home are better. The beds are more comfortable. It's more peaceful. And at home, the food is better." I'll admit she's right about the first three. But the food? We had a simple pasta one night, snags another, chicken curry the third, but they all tasted good. Better than at home. And it wasn't simply that we were hungry; it was because we were camping.\par
\par The real point I want to get to is that meals largely taste so good when I'm camping because I'm reminded of some things easily forgotten at home, where I can simply go to the cupboard, or the supermarket, and get whatever takes my fancy. Life, when you're camping, is reduced to its essentials, and I am reminded of how lucky most of us are to be dry and warm and fed. More importantly, I am reminded that sadly, and in fact unforgivably, not all of us are. And I am reminded that we, as a society, should do better for them.\par
\question{24. Why did Jared bring only rice for dinner on the trip?}
\begin{choices}(2)
  \item He didn't like soy sauce.
  \item He preferred to travel light.
  \item He couldn't afford a better meal.
  \item He enjoyed simple outdoor meals.
\end{choices}
\question{25. What does the underlined word "ravenous" in the third paragraph mean?}
\begin{choices}(2)
  \item Tired.
  \item Starving.
  \item Fulfilled.
  \item Refreshed.
\end{choices}
\question{26. What does the author think of camping?}
\begin{choices}(2)
  \item It is a test of adaptability.
  \item It is a chance to enjoy good food.
  \item It is a reminder of the good life he has.
  \item It is a bonding experience for his family.
\end{choices}
\question{27. What makes the author sad according to the last paragraph?}
\begin{choices}(2)
  \item Some people lack the basic necessities of life.
  \item Some people fail to appreciate what they have.
  \item People easily forget the beauty of simple things.
  \item People easily quit camping due to its discomfort.
\end{choices}
\passagelabel{C}
\par Neighboring groups of birds within the same species can create very different-looking nests — showing that their nest-building choices aren't only controlled by instinct(本能) and the environment. Instead, these birds seem to learn rules for nest-making that get passed down within a family group from generation to generation.\par
\par That's the conclusion of researchers who studied nearly 450 grassy structures built over two years by white-browed sparrow weavers living in the Kalahari Desert in South Africa. These small, brown-and-white birds live in groups, and it turns out that groups of birds that lived quite near to each other built distinctive architectural forms.\par
\par The biggest difference was in "how short or long the structures are," says study lead author Maria C. Tello-Ramos, a cognitive ecologist at the University of Hull in England. Furthermore, each group maintained the same architectural style over time—and when outsiders joined, they adapted to this style.\par
\par To examine what might account for the observed differences, the team carefully analyzed factors that can determine a nest's size and shape for a given bird species: weather conditions, tree height, individuals' body size and genetic(基因的) relatedness. Yet none of these factors seemed to play a relevant role in shaping how the Kalahari sparrows built their nests, the researchers report in Science.\par
\par "Then we say, ‘Okay, so what is left?’” Tello-Ramos explains. She and her colleagues proposed that cultural transmission(传输) might be key to nest building. "In our paper, we haven’t gotten there yet with experiments, but we have very good clues that that might be it," she says.\par
\par "These are important questions that are understudied," says Christina Riehl, an evolutionary biologist. She's not convinced the study's data are enough to fully rule out genetic influence. "They can't actually look at the effect of, say, genetic differences, because they don't have really good genetic information on all the individuals in these groups," she says. "I think there's a lot left to be done, and I think this paper will inspire future research in a really good way."\par
\question{28. What plays a key role in birds' nest-building choices according to the new study?}
\begin{choices}(2)
  \item Living conditions.
  \item Group traditions.
  \item The influence of neighbors.
  \item The species of birds.
\end{choices}
\question{29. What do we know about the nests built by neighboring groups of birds?}
\begin{choices}(2)
  \item They are modeled after each other.
  \item They are made of different materials.
  \item They vary most significantly in length.
  \item They share the same architectural style.
\end{choices}
\question{30. How did Tello-Ramos and his colleagues reach the study conclusion?}
\begin{choices}(2)
  \item Through statistical methods.
  \item Through careful experiments.
  \item By excluding other possible factors.
  \item By reviewing findings of previous studies.
\end{choices}
\question{31. What does Christina Riehl think of the study?}
\begin{choices}(2)
  \item Limited but inspiring.
  \item Understudied and misleading.
  \item Controversial but well-designed.
  \item Convincing and groundbreaking
\end{choices}
\passagelabel{D}
\par A wheeled robot rolls across the floor. A soft-bodied robotic star bends its five legs, moving awkwardly. Powered by conventional electricity via plug or battery, these simple robotic creations would be unremarkable, but what sets these two robots apart is that they are controlled by a living entity(实体): a king oyster mushroom.\par
\par By growing the mushroom's mycelium(菌丝体) into the robot's hardware, a team led by Cornell University researchers has engineered two types of robots that sense and respond to the environment by using electrical signals made by the mushroom and its sensitivity to light.\par
\par The robots are the latest accomplishment of scientists in a field known as biohybrid robotics who seek to combine biological, living materials such as plant and animal cells or insects with artificial components to make partly living and partly engineered entities\par
\par There's still a long way to go before biohybrid robots go beyond the lab, but researchers hope one day robot jellyfish may explore oceans, sperm-powered robots may be able to deliver fertility(肥沃) treatments and cyborg cockroaches could search for survivors in the wake of an earthquake.\par
\par "Biohybridization is an attempt to find components in the biological world that we can use, understand, and control to help our artificial systems work better," said Shepherd, a professor of mechanical and aerospace engineering at Cornell University.\par
\par The team began by growing king oyster mushrooms in the lab from a simple kit ordered online. The researchers chose this species of mushroom because it grows easily and quickly. They cultivated(培育) the mushroom's mycelium, which can form networks that, according to the study, can sense, communicate and transport nutrients.\par
\par The team found it challenging to engineer a system that could detect and use the small electrical signals from the mycelia to command the robot. "You have to make sure that your electrode(电极) touches in the right position because the mycelia are very thin," said lead author Anand Mishra. "Then you culture them, and when the mycelia start growing, they wrap around the electrode."\par
\question{32. What is special about the two robots in the first paragraph?}
\begin{choices}(2)
  \item They can control living creatures
  \item They can make various movements.
  \item They are powered by king oyster mushrooms.
  \item They are partly mushroom and partly machine.
\end{choices}
\question{33. What does the author want to show by giving the examples in paragraph 4?}
\begin{choices}(2)
  \item The future possibilities of biohybrid robots.
  \item The seriousness of the environmental issues.
  \item The limitations of biohybrid robots beyond the lab.
  \item The advantages of biohybrid robots over conventional ones.
\end{choices}
\question{34. What did the team find difficult when engineering the robots?}
\begin{choices}(2)
  \item Wrapping the electrode with the mycelia.
  \item Arranging the electrode in the right order.
  \item Developing a system to cultivate the mycelia.
  \item Positioning the electrode accurately on the mycelia
\end{choices}
\question{35. Which of the following can be a suitable title for the text?}
\begin{choices}(2)
  \item When Nature Meets Technology
  \item How Technology Transforms Nature
  \item King Oyster Mushrooms: the Magic of Nature
  \item Biohybrid Robotics: the Key to Social Development
\end{choices}
\sectitle{第二节（共5小题；每小题2.5分，满分12.5分）}
\par 阅读下面短文，从短文后的选项中选出能填入空白处的最佳选项。选项中有两项为多余选项，\par
\par Science shows that attraction isn't just about looks. \blanknum{36} The following simple ways will help you boost attractiveness and impact how others perceive you.\par
\par Master the power of a smile. A genuine smile can significantly boost your attractiveness. Research shows that a better smile enhances facial appeal for both men and women. Smiling not only makes you look more approachable but also signals warmth, confidence, and positive emotions. \blanknum{37}\par
\par Be a good audience. Want to increase your appeal instantly? Laugh at someone's jokes sincerely. A study conducted in the United States and Norway found that expressing amusement — giggling or laughing — when someone makes a joke significantly increases romantic interest. \blanknum{38} As a result, interactions become more engaging and attractive.\par
\par \blanknum{39}Appearing energetic and lively boosts attractiveness. Facial expressions and body language that convey liveliness, activity, and enthusiasm make someone seem more appealing. Since videos capture liveliness better than static(静态的) images, if you're using dating apps or social media, consider using videos instead of just photos to highlight your dynamic presence.\par
\par Maximize the "cheerleader effect". Want to appear more attractive? Make the best use of the "cheerleader effect" by spending time with a group of people that includes at least some attractive individuals. \blanknum{40} So, next time you want to make a great impression, gather your charming friends and let the magic of the group work for you!\par
\par Attractiveness is about embracing your best qualities, building confidence, and forming meaningful connections. Remember, the most attractive version of you is the one that feels happy, empowered, and unapologetically yourself.\par
\begin{choices}(1)
  \item Be real in videos.
  \item Make your energy obvious.
  \item All of them are highly attractive qualities.
  \item Surround yourself with a group of interesting people
  \item Humor creates a sense of connection and playfulness.
  \item It's about how you carry yourself, connect with others, and highlight your best qualities.
  \item It's because the brain averages facial features in groups, unnoticeably improving perceived beauty.
\end{choices}
\parttitle{第三部分 语言知识运用（共两节，满分 30 分）}
\sectitle{第一节（共15小题；每小题1分，满分15分）}
\par 阅读下面短文，从短文后各题所给的 A、B、C 和 D 四个选项中，选出可以填入空白处的最佳选项，\par
\par My late mother was a famous artist in the Pacific Northwest. Over her many-decade career, her paintings kept changing. One \blanknum{41} in her work, however, was excellent technique: If she \blanknum{42} to paint a naked guy holding a guitar, that's \blanknum{43} what it looked like.\par
\par Growing up, I could draw a little myself, but I never had her \blanknum{44}. Once, I asked her how I could improve. I \blanknum{45} her to say something like "Practice 10,000 hours." Instead, she told me to \blanknum{46} what I wanted to draw. This \blanknum{47} me because that's obviously what I thought I had been doing.\par
\par "People almost never actually look carefully at anything," she \blanknum{48}, "They glance(臀) at it and then rely on their brain to fill in the \blanknum{49} – which it doesn’t, leading to \blanknum{50} drawing." So I tried, looking long and hard at what I wanted to draw: a tree. I found I \blanknum{51} much more about its shape, colors, and shadows. I \blanknum{52} each detail carefully, and it turned out to be a pretty well-drawn tree.\par
\par More than that, I loved the experience of really \blanknum{53}. What my mom was telling me to do was fully enjoy the encounter of seeing something \blanknum{54} and drawing it carefully. And it can be \blanknum{55} to many areas of life in ways that help us become better at living.\par
\cz{41}{favorite}{constant}{potential}{alternative}
\cz{42}{preferred}{pretended}{decided}{struggled}
\cz{43}{exactly}{probably}{naturally}{strangely}
\cz{44}{skill}{permission}{patience}{fame}
\cz{45}{urged}{reminded}{encouraged}{expected}
\cz{46}{determine}{picture}{observe}{introduce}
\cz{47}{delighted}{relieved}{angered}{confused}
\cz{48}{apologized}{explained}{announced}{repeated}
\cz{49}{colors}{details}{outline}{pattern}
\cz{50}{awkward}{abstract}{original}{attractive}
\cz{51}{accepted}{doubted}{mentioned}{noticed}
\cz{52}{analyzed}{remembered}{described}{drew}
\cz{53}{thinking}{sharing}{looking}{improving}
\cz{54}{deeply}{frequently}{directly}{differently}
\cz{55}{exposed}{compared}{applied}{limited}
\sectitle{第二节（共10小题；每小题1.5ladnepen分，满分15分）}
\par 阅读下面短文，在空白处填入 1 个适当的单词或括号内单词的正确形式。\par
\par Rows of delicate liuli (colored glass), from vases to teacups, shine in the light at SunYunyi's studio in Zibo, Shandong province. Among them, the most striking pieces are in the color that looks like the fat under a hen's stomach, and thus are known \blanknum{56} jiyouhuang (鸡油黄).\par
\par Jiyouhuang came into being during the Ming Dynasty and were mostly produced during the Qing Dynasty. Over time, research related to its production and \blanknum{57} (survive) artifacts(手工制品) have been lost, making it \blanknum{58} increasingly rare historical treasure.\par
\par Using bits of information from his ancestors, \blanknum{59} used to work in the royal workshop, Sun and his family managed \blanknum{60} (rediscover) the way jiyouhuang was made in 2007.\par
\par There was no existing documentation of the full recipe and \blanknum{61} (compose) of jiyouhuang, so initially, the success rate was low. Sun explored the use of other raw materials to adjust the recipe \blanknum{62} improved craftsmanship(技艺) accordingly to make the new recipe better for iyouhuang, giving rise to a much \blanknum{63} (high) success rate and purity of product.\par
\par According to evaluation by the Palace Museum in Beijing, his jiyouhuang pieces are better than the Qing Dynasty \blanknum{64} (original) in quality. In 2013, Sun became a national patent(专利) holder for the jiyouhuang production method, and the craft \blanknum{65} (name) a provincial intangible cultural heritage in Shandong in 2016.\par
\parttitle{第四部分 写作（共两节，满分 40 分）}
\sectitle{第一节（满分15分）}
\par 假定你是某国际学校学生李华。为了帮助师生放松身心，提升校园幸福感，你校打算设立一个“元气补给站”(A Refresh Corner)，现向全体师生征求设计意见，请你写一篇短文，内容包括：\par
\par (1) 描述你的设计；\par
\par (2) 阐述设计理由。 注意：\par
\par (1) 写作词数应为 80 左右;\par
\par (2) 请按如下格式在答题纸的相应位置作答。\par
\par (2) 请按如下格式在答题纸的相应位置作答。 A Refresh Corner on Campus\par
\sectitle{第二节（满分25分）}
\par 阅读下面材料，根据其内容和所给段落开头语续写两段，使之构成一篇完整的短文。\par
\par In the neighbourhood, my dad, Mr. Duane, is always referred to as "Collector Duane". He is an enthusiastic collector, who is always on the hunt for old items to repurpose, believing that with a bit of creativity, these forgotten objects can be transformed into something useful and new. Recently, he came across an old bicycle wheel, an old-fashioned metal milk box, an antique lamp, a dirty vase and some old clothes at a flea market, imagining all the possibilities they held.\par
\par My parents have a neighbor: Helen Fox. She is a 91-year-old widow(寡妇). My parents, along with several other neighbors, look after Helen, ensuring she feels loved and supported.\par
\par Helen had long suffered from a problem: her mailbox. For some reason, the mailbox was placed insecurely down a somewhat steep(陡峭的) driveway off a busy highway, making it a daily danger for her to fetch her mail. Besides, it was old and rusty, as if it had suffered multiple violent impacts. Its once-bright red paint had faded to a dull pink. The door hung loosely, creaking(嘎吱作响) with every gust of wind. The letters on it had faded over the years. The 'F' was no longer clear, the O' had lost its shape, and the 'X' was almost gone. Helen wanted to replace it with a new one and move it to a new location, but she had no idea where to start.\par
\par The minute my dad heard about it, he made it his mission to settle the problem. He began by contacting the local rural mail service, arranging to relocate Helen's mailbox. He also talked with Yancey, their shared rural postman, explaining the upcoming change. A relatively safe place was easy to find. But the real challenge was finding a perfect mailbox.\par
\par 注意:\par
\par (1) 续写词数应为 150 左右;\par
\par (2) 请按如下格式在答题卡的相应位置作答。\par
\writingpara{Paragraph 1: Then he remembered the old objects he had recently collected. \blanknum{}}
\writinglines{8}
\writingpara{Paragraph 2: The next day, he showed up at Helen's house with a shiny new mailbox. \blanknum{}}
\writinglines{8}
\end{document}
